% XeLaTeX 编译: xelatex collaboration_proposal_zhang_zh.tex
\documentclass[a4paper,11pt]{article}
\usepackage[a4paper, top=2.5cm, bottom=2.5cm, left=2.5cm, right=2.5cm]{geometry}
\usepackage{fontspec}
\usepackage[UTF8]{ctex}
\usepackage{microtype}
\usepackage{amsmath,amssymb}
\usepackage{booktabs}
\usepackage{array}
\usepackage{hyperref}
\usepackage{graphicx}
\usepackage{parskip}
\usepackage{fancyhdr}
\usepackage{titlesec}

\hypersetup{
  pdftitle={合作提案 RFT -- 张靖教授},
  pdfauthor={Dominic Ren\'e Schu},
  colorlinks=true,
  linkcolor=black,
  urlcolor=blue,
  citecolor=black
}

\pagestyle{fancy}
\fancyhf{}
\lhead{\small 共振场理论 --- 合作提案}
\rhead{\small Dominic Ren\'e Schu}
\cfoot{\thepage}
\renewcommand{\headrulewidth}{0.4pt}

\titleformat{\section}{\large\bfseries}{}{0em}{}[\titlerule]
\titleformat{\subsection}{\normalsize\bfseries}{}{0em}{}

\setlength{\parindent}{0pt}

%% ====================================================================
\begin{document}

%% ---- 封面 ----
\thispagestyle{empty}
\begin{center}
  {\LARGE\bfseries 合作提案}\\[0.5em]
  {\large\bfseries $^{87}$Rb 玻色-爱因斯坦凝聚体中\\
  非线性量子力学的可证伪预测}\\[2em]
  {\large Dominic Ren\'e Schu}\\[0.3em]
  {\normalsize 共振场理论（RFT）\\
  \href{https://github.com/DominicReneSchu/RFT}{github.com/DominicReneSchu/RFT}}\\[2em]
  \rule{0.6\textwidth}{0.4pt}\\[1em]
  {\normalsize 致：}\\[0.3em]
  {\large\bfseries 张靖 教授}\\
  {\normalsize 量子光学与光学装置国家重点实验室\\
  山西大学，太原，中国}\\[2em]
  {\normalsize \today}
\end{center}

\vfill
\begin{center}
\small
本文件为正式科学合作提案。\\
完整理论基础、所有模拟程序及实验方案均已公开发布于：\\
\href{https://github.com/DominicReneSchu/RFT}{https://github.com/DominicReneSchu/RFT}
\end{center}
\newpage

%% ---- 摘要 ----
\section{摘要}

\textbf{共振场理论}（Resonance Field Theory，RFT）是一个公理化框架，将薛定谔方程推导为
更一般非线性动力学的特殊情形。通过密度反馈模型
\[
  \dot{\Delta\varphi} = \lambda \int |\psi|^4\,\mathrm{d}x,
\]
RFT 在极限 $\lambda \to 0$ 时精确还原为标准量子力学，而在 $\lambda > 0$ 时给出可测量的偏差。

本合作提案提出一个完整的\textbf{可证伪实验方案}：
\[
  \boxed{|\Delta\langle x\rangle|_{\mathrm{SI}} \approx 2{.}0 \cdot \lambda\;\mu\mathrm{m}}
\]
该预测基于数值验证的标度律、针对 $^{87}$Rb 谐振阱的完整 SI 标定
以及严格表述的统计检验协议。

\textbf{合作目标：}利用张靖教授课题组（山西大学）现有的 BEC 实验装置
对上述预测进行实验检验。该实验为\emph{界限实验}：
要么检测到 RFT 效应，要么为 $\lambda$ 设定上限——
两种结果均具有重要科学价值。

\newpage
\tableofcontents
\newpage

%% ====================================================================
\section{理论背景}

\subsection{共振场理论：公理四与薛定谔方程}

RFT 建立在七条公理之上。对于本实验而言，\textbf{公理四}（耦合效率）最为关键：
量子系统与共振场之间的有效耦合由相位差 $\Delta\varphi$ 调制：
\[
  \varepsilon(\Delta\varphi) = \cos^2\!\left(\frac{\Delta\varphi}{2}\right) \in [0, 1].
\]
共振哈密顿算符为：
\[
  \hat{H}_{\mathrm{res}} = \hat{H}_0 + \varepsilon(\Delta\varphi)\,\hat{V}_{\mathrm{耦合}}.
\]
对于恒定的 $\Delta\varphi$（静态极限），可精确还原为标准薛定谔方程——
对应原理已通过解析和数值方法证实（所有测试 $\Delta\varphi$ 对应的保真度均为
$F = 1{.}000000000000$）。

\subsection{密度反馈模型}

在动力学推广中，$\Delta\varphi(t)$ 成为与波函数反馈耦合的时间相关场：
\[
  \Delta\varphi(t + \mathrm{d}t) = \Delta\varphi(t)
  + \lambda \cdot \underbrace{\int |\psi(x,t)|^4\,\mathrm{d}x}_{F[\psi]} \cdot \mathrm{d}t.
\]
该模型可从 RFT 作用量泛函通过欧拉-拉格朗日方程推导（过阻尼极限）：
\[
  S[\psi, \Delta\varphi] = \int \mathrm{d}t \left[
    \langle\psi|i\hbar\partial_t - \hat{H}_0|\psi\rangle
    - \varepsilon(\Delta\varphi)\langle V\rangle_\psi
    + \frac{\mu}{2}(\dot{\Delta\varphi})^2
  \right].
\]
反馈链 $\psi \to \Delta\varphi \to V_{\mathrm{eff}} \to \psi$ 严格局域，
不违反不可信号传输原理（No-Signaling）。

\subsection{微扰理论与标度律}

对于 $\lambda \ll 1$，一阶微扰理论给出所有可观测量的解析表达式。
数值计算以相对误差 $1{.}3 \times 10^{-4}$ 验证了这些结果。
核心标度律如下：

\begin{table}[h]
\centering
\caption{RFT 可观测量的标度指数（微扰理论与数值计算的比较）}
\begin{tabular}{lccc}
\toprule
可观测量 & 指数（数值） & 指数（理论） & 偏差 \\
\midrule
$1 - F$（保真度） & $2{.}001$ & $2$ & $0{.}05\,\%$ \\
$|\Delta\langle x\rangle|$ & $1{.}001$ & $1$ & $0{.}1\,\%$ \\
$|\Delta\langle p\rangle|$ & $1{.}001$ & $1$ & $0{.}1\,\%$ \\
$\max|\Delta\psi|$ & $0{.}999$ & $1$ & $0{.}1\,\%$ \\
\bottomrule
\end{tabular}
\end{table}

数值确定的前置因子为：
\[
  |\Delta\langle x\rangle| \approx 4{.}9 \cdot \lambda, \quad
  |\Delta\langle p\rangle| \approx 3{.}2 \cdot \lambda, \quad
  1 - F \approx 138 \cdot \lambda^2.
\]
$1 - F \sim \lambda^2$ 的关系直接源于
$|\langle\psi_0|\psi_0 + \lambda\psi_1\rangle|^2
= 1 - \lambda^2\|\psi_1\|^2 + \mathcal{O}(\lambda^3)$。

%% ====================================================================
\section{SI 单位标定}

\subsection{物理系统}

参考系统为\textbf{谐振阱中的超冷 $^{87}$Rb 原子}——
全球精密量子力学实验中应用最广泛的标准系统：

\begin{table}[h]
\centering
\caption{$^{87}$Rb 系统参数}
\begin{tabular}{llr}
\toprule
参数 & 符号 & 数值 \\
\midrule
原子质量 & $m$ & $86{.}909\,\mathrm{u} = 1{.}443 \times 10^{-25}\,\mathrm{kg}$ \\
阱频率 & $\omega$ & $2\pi \times 100\,\mathrm{Hz}$（可调） \\
约化普朗克常数 & $\hbar$ & $1{.}055 \times 10^{-34}\,\mathrm{J\cdot s}$ \\
\bottomrule
\end{tabular}
\end{table}

\subsection{完整单位换算}

无量纲模拟（$\hbar = m = 1$）通过三个转换量映射到 SI 单位。
标定由 6 个独立的一致性关系（如 $\hbar = m\,\ell^2/\tau$，
$E \cdot \tau = \hbar$）过定，完全自洽。

\begin{table}[h]
\centering
\caption{$\omega = 2\pi \times 100\,\mathrm{Hz}$ 时的 SI 标定}
\begin{tabular}{llll}
\toprule
物理量 & 公式 & 数值（$\omega = 2\pi \times 100\,\mathrm{Hz}$） \\
\midrule
谐振子长度 & $a_{\mathrm{ho}} = \sqrt{\hbar/(m\omega)}$ & $1{.}08\,\mu\mathrm{m}$ \\
长度单位 & $\ell = V_s^{1/4} \cdot a_{\mathrm{ho}}$ & $0{.}41\,\mu\mathrm{m}$ \\
时间单位 & $\tau = \sqrt{V_s}/\omega$ & $0{.}225\,\mathrm{ms}$ \\
能量单位 & $E = \hbar/\tau$ & $2{.}92 \times 10^{-12}\,\mathrm{eV}$ \\
动量单位 & $p = \hbar/\ell$ & $2{.}60 \times 10^{-28}\,\mathrm{kg\cdot m/s}$ \\
模拟时长 & $T = 20\,\tau$ & $4{.}5\,\mathrm{ms}$ \\
\bottomrule
\end{tabular}
\end{table}

\subsection{单位换算推导}

无量纲薛定谔方程（$\hbar = m = 1$）：
\[
  i\,\partial_{\tilde{t}}\,\psi = -\tfrac{1}{2}\,\partial_{\tilde{x}}^2\,\psi
  + V(\tilde{x})\,\psi
\]
通过 $x = \ell\,\tilde{x}$，$t = \tau\,\tilde{t}$ 映射到 SI 单位。
谐振耦合势 $V(\tilde{x}) = \tfrac{1}{2}V_s\,\tilde{x}^2$
等同于物理势 $V_{\mathrm{phys}}(x) = \tfrac{1}{2}m\omega^2 x^2$：
\[
  V_s = \frac{m^2\omega^2\ell^4}{\hbar^2}
  \quad\Longrightarrow\quad
  \ell = V_s^{1/4} \cdot \sqrt{\frac{\hbar}{m\omega}} = V_s^{1/4} \cdot a_{\mathrm{ho}}.
\]

%% ====================================================================
\section{可证伪预测}

\subsection{核心预测}

RFT 预测波包质心位移：
\[
  \boxed{|\Delta\langle x\rangle|_{\mathrm{SI}}
  = C_x \cdot \lambda \cdot \ell
  \approx 2{.}0 \cdot \lambda\;\mu\mathrm{m}}
\]
其中 $C_x \approx 4{.}9$（数值确定），
$\ell \approx 0{.}41\,\mu\mathrm{m}$（$\omega = 2\pi \times 100\,\mathrm{Hz}$ 时）。

\subsection{预测数值表}

\begin{table}[h]
\centering
\caption{不同 $\lambda$ 值对应的位置偏移和保真度损失预测}
\begin{tabular}{rrrrl}
\toprule
$\lambda$ & $|\Delta\langle x\rangle|$ [$\mu$m] & $1 - F$ & 可探测性 \\
\midrule
$0{.}001$ & $0{.}002$ & $1{.}4 \times 10^{-4}$ & 否（单次测量） \\
$0{.}01$  & $0{.}020$ & $0{.}014$               & 是（$10^4$ 次重复） \\
$0{.}05$  & $0{.}099$ & $0{.}345$               & 是（100 次重复） \\
$0{.}1$   & $0{.}199$ & $1{.}38$                & 是（100 次重复） \\
$0{.}5$   & $0{.}994$ & $34{.}5$                & 是（单次测量） \\
$1{.}0$   & $1{.}987$ & $138$                   & 是（单次测量） \\
\bottomrule
\end{tabular}
\end{table}

\subsection{阱频率依赖性与最优区间}

较低的阱频率对应更大的 $a_{\mathrm{ho}}$，从而提升探测灵敏度：

\begin{table}[h]
\centering
\caption{灵敏度随阱频率的变化}
\begin{tabular}{rrrr}
\toprule
$\omega / 2\pi$ [Hz] & $a_{\mathrm{ho}}$ [$\mu$m] & $\ell$ [$\mu$m]
  & $\lambda_{\min}$（100 次重复） \\
\midrule
$10$   & $3{.}41$ & $1{.}28$ & $0{.}016$ \\
$50$   & $1{.}53$ & $0{.}57$ & $0{.}036$ \\
$100$  & $1{.}08$ & $0{.}41$ & $0{.}050$ \\
$500$  & $0{.}48$ & $0{.}18$ & $0{.}113$ \\
$1000$ & $0{.}34$ & $0{.}13$ & $0{.}159$ \\
\bottomrule
\end{tabular}
\end{table}

\textbf{最优区间：} $\omega \lesssim 2\pi \times 50\,\mathrm{Hz}$
（但制备时间较长）。

%% ====================================================================
\section{实验方案}

\subsection{制备阶段}

\begin{enumerate}
  \item \textbf{BEC 制备：} 在磁阱或光阱中制备 $^{87}$Rb 凝聚体，
    $T < 200\,\mathrm{nK}$，$N \sim 10^5$ 个原子。
  \item \textbf{谐振阱：} 轴向频率可调至
    $\omega = 2\pi \times (10\text{--}500)\,\mathrm{Hz}$。
    轴向频率决定 $a_{\mathrm{ho}}$ 进而决定探测灵敏度。
  \item \textbf{波包初始化：} 通过布拉格脉冲或拉曼跃迁施加动量冲激：
    $\hbar k_0 = \hbar/\ell$。
\end{enumerate}

\subsection{测量阶段}

\begin{enumerate}
  \setcounter{enumi}{3}
  \item \textbf{自由演化：} 波包在阱中演化
    $t \approx 4{.}5\,\mathrm{ms}$（$\omega = 2\pi \times 100\,\mathrm{Hz}$ 时）。
  \item \textbf{吸收成像：} 关闭阱，飞行时间展开，
    采集吸收图像。空间分辨率：${\sim}\,1\,\mu\mathrm{m}$。
  \item \textbf{重复测量：} $N = 100\text{--}10\,000$ 次相同实验序列。
\end{enumerate}

\subsection{数据分析}

\begin{enumerate}
  \setcounter{enumi}{6}
  \item \textbf{统计分析：} 从 $N$ 次测量中提取平均位置
    $\langle x\rangle_{\mathrm{exp}}$。
    标准误：$\sigma_{\Delta x} = \sigma_{\mathrm{single}}/\sqrt{N}$。
  \item \textbf{假设检验：}
    \begin{itemize}
      \item \textbf{零假设} $H_0$：$\Delta\langle x\rangle = 0$
        （标准量子力学，$\lambda = 0$）
      \item \textbf{备择假设} $H_1$：
        $|\Delta\langle x\rangle| = C_x \cdot \lambda \cdot \ell > 0$（RFT）
    \end{itemize}
  \item \textbf{实验结果：}
    \begin{itemize}
      \item $|\Delta\langle x\rangle| > 2\sigma$：探测到 RFT 效应
        $\to$ 定量确定 $\lambda$
      \item $|\Delta\langle x\rangle| \leq 2\sigma$：设定上限
        $\lambda_{\max}$
    \end{itemize}
\end{enumerate}

%% ====================================================================
\section{为何选择张靖教授课题组}

张靖教授在山西大学量子光学与光学装置国家重点实验室领导的课题组
恰好具备本实验方案所需的全部实验设施和专业技术：

\begin{itemize}
  \item \textbf{$^{87}$Rb BEC 专业积累：}
    在铷玻色-爱因斯坦凝聚体的制备与操控方面拥有多年经验——
    这正是本方案所要求的精确实验系统。
  \item \textbf{吸收成像技术：}
    空间分辨率 ${\sim}\,1\,\mu\mathrm{m}$ 的高分辨吸收成像
    是本方案的核心测量手段。
  \item \textbf{相干动力学与波包操控：}
    已发表的相干波包动力学工作与本方案的测量需求
    （动量冲激、自由演化、位置测量）高度匹配。
  \item \textbf{国家重点实验室基础设施：}
    磁场控制、可调谐阱频率与光学偶极阱
    均在方案中明确规划使用。
  \item \textbf{与实验方案的直接对应：}
    BEC 制备、谐振阱操控、动量冲激和吸收成像
    均是张靖课题组的常规实验技术。
\end{itemize}

%% ====================================================================
\section{批判性评估}

\subsection{Kohn 定理：与 Gross-Pitaevskii 效应的区分}

RFT 反馈项 $\dot{\Delta\varphi} \propto \int |\psi|^4\,\mathrm{d}x$
在形式上与 Gross-Pitaevskii 方程中的接触相互作用项（$g|\psi|^2\psi$）相同。
概念上的区分由 \textbf{Kohn 定理}给出：

在\emph{纯谐振阱}中，质心运动 $\langle x\rangle(t)$ 精确地以频率 $\omega$
做谐振——\emph{与原子间相互作用无关}：
\[
  \text{GP：} \quad \Delta\langle x\rangle = 0 \quad\text{（Kohn 保护）}.
\]
RFT 反馈则通过 $\varepsilon(\Delta\varphi(t)) \cdot V$ 对阱强度进行时间依赖性调制，
从而\emph{破坏} Kohn 条件：
\[
  \text{RFT：} \quad |\Delta\langle x\rangle| = C_x \cdot \lambda \cdot \ell \neq 0.
\]
GP 项改变凝聚体的\emph{宽度}，而非\emph{质心}。
因此 RFT 信号在概念上与平均场效应清晰分离。

\subsection{系统误差分析}

\begin{table}[h]
\centering
\caption{$^{87}$Rb BEC 谐振阱实验的误差预算}
\begin{tabular}{lrl}
\toprule
误差来源 & 估计位移 & 标度关系 \\
\midrule
GP 平均场对 $\langle x\rangle$ 的贡献 & \textbf{0}（Kohn 定理） & 不适用 \\
势阱非简谐性 & ${\sim}\,1\,\mathrm{nm}$ & $\delta\omega/\omega$ \\
磁场梯度（$0{.}1\,\mathrm{mG/cm}$） & ${\sim}\,650\,\mathrm{nm}$ & $\mu_B \cdot \nabla B$ \\
三体损耗 & ${\sim}\,11\,\mathrm{nm}$ & $\Delta N / N$ \\
\midrule
\textbf{总计（平方叠加）} & $\mathbf{{\sim}\,650\,\mathrm{nm}}$ & \\
\bottomrule
\end{tabular}
\end{table}

对于 $\lambda \lesssim 0{.}3$，磁场梯度为主导误差来源。
将磁场补偿优化至 $< 0{.}01\,\mathrm{mG/cm}$ 可使系统误差降低一个量级。

\subsection{区分实验方案}

\begin{table}[h]
\centering
\caption{GP 与 RFT 的实验区分方案}
\begin{tabular}{lll}
\toprule
扫描方案 & GP 依赖性 & RFT 依赖性 \\
\midrule
$N$ 扫描 & 宽度 $\propto N^{2/5}$ & $\langle x\rangle$ 偏移 $\neq f(N)$ \\
$a_s$ 扫描（Feshbach） & $\propto a_s$ & 与 $a_s$ 无关 \\
$\omega$ 扫描 & 宽度 $\propto \omega^{-3/5}$ & 偏移 $\propto \omega^{-1/2}$ \\
$\Delta\varphi_0$ 扫描 & 无关 & $\propto \varepsilon'(\Delta\varphi_0)$ \\
\bottomrule
\end{tabular}
\end{table}

\subsection{作为界限实验的定位}

本实验应透明地定位为\emph{界限实验}：
\begin{itemize}
  \item \textbf{正结果：} $|\Delta\langle x\rangle| > 2\sigma$
    $\Rightarrow$ 探测到 RFT 效应，定量确定 $\lambda$。
  \item \textbf{零结果：} $|\Delta\langle x\rangle| \leq 2\sigma$
    $\Rightarrow$ 设定上限 $\lambda_{\max}$，对 RFT 参数空间形成显著约束。
\end{itemize}
两种实验结果均具有科学价值，均可发表于高水平同行评审期刊。

%% ====================================================================
\section{合作框架建议}

\begin{itemize}
  \item \textbf{理论与模拟：}
    Dominic Ren\'e Schu——完整解析基础、数值模拟、
    SI 标定、预测计算。
  \item \textbf{实验：}
    张靖课题组（山西大学）——$^{87}$Rb BEC 制备、
    谐振阱操控、吸收成像、统计分析。
  \item \textbf{联合发表：}
    双方作为共同作者在同行评审期刊发表
    （如 \textit{Physical Review Letters} 或 \textit{Nature Physics}）。
  \item \textbf{开放科学：}
    所有模拟程序、数据和分析结果
    在 GitHub 仓库公开记录。
\end{itemize}

%% ====================================================================
\section{参考资料}

\begin{enumerate}
  \item Dominic Ren\'e Schu，\textit{共振场理论——公理化框架与实验方案}，
    GitHub 仓库，
    \url{https://github.com/DominicReneSchu/RFT}（2025/2026）。
  \item 完整实验方案（SI 标定、微扰理论、可证伪预测）：\\
    \texttt{de/fakten/simulationen/schrödinger/docs/experimental\_proposal.md}
  \item 模拟程序包（Python，numpy）：\\
    \texttt{de/fakten/simulationen/schrödinger/python/schrodinger\_1d\_rft\_experiment.py}
  \item 薛定谔模拟概述：\\
    \texttt{de/fakten/simulationen/schrödinger/README.md}
  \item W.~Kohn，\textit{Cyclotron Resonance and de Haas-van Alphen
    Oscillations of an Interacting Electron Gas}，
    Phys.\ Rev.\ \textbf{123}，1242（1961）。\,[\,Kohn 定理\,]
  \item C.~Pethick 与 H.~Smith，
    \textit{Bose-Einstein Condensation in Dilute Gases}，
    剑桥大学出版社（2008）。
\end{enumerate}

\vspace{2em}
\hrule
\vspace{1em}
\begin{center}
\small
\textit{Dominic Ren\'e Schu --- 共振场理论}\\
\href{https://github.com/DominicReneSchu/RFT}{https://github.com/DominicReneSchu/RFT}
\end{center}

\end{document}
